\scanpage{49}
\mspage{7}
rencontre la diagonale, OK.

\uline{Corollaire 3} Soit $K'$ une \uline{extension} \uline{primaire} d'un corps $k$, $G$ un schéma en groupes comm de type fini sur $k$, $H'$ un sous schéma \inserted{en groupes fermé} de $G'=G\otimes_k k'$, \uline{rigide}. Alors il existe un ss-schéma \inserted{en groupes} fermé $H$ de $G$, unique, \deleted{rigide,} tel que $H'=H\otimes_k k'$.

On admettra qu'il existe une sous-algèbre \deleted{\unsure{un} $\unsure{K_1|A_1}$}\inserted{$A_1$} de $K'$ \inserted{de} type fini sur $k$, telle que $H'$ provient d'un ss schéma en groupes fermé $H'_1$ de $G_1\deleted{\unsure{'}}=G\otimes_{\unsure{A|k}}A_1^{\deleted{\unsure{'|*}}}$, \deleted{\unsure{si}} qui \uline{est rigide}

% slanted
\margin{préciser ce point}
[c'est le cas en tous cas si on \deleted{\unsure{admet}} suppose que $H'$ a, sur la clôture algébrique de $K'$, une suite de composition dont les facteurs sont des $G_m$, $Ab$, $\mathbb Z/\ell\mathbb Z\ \unsure{\text{---}|\ldots}$].

\deleted{On peut alors \unsure{On}} Comme le corps $K_1$ de $A_1$ est une extension primaire de $k$, on sait que $\unsure{\operatorname{hom.}|\operatorname{hom}|\operatorname{Spec}}(A_1)$ est \inserted{\unsure{abt}} connexe. On applique alors le cor. 2.
% archive sheet number: 48
